\chapter{Odometry}
\label{ch:odometry}

\pkg{dharma_odometry} starts the LiDAR-inertial estimator behind one launch file, holds its
start until the IMU is still and level, and monitors its output. The estimator publishes
\frm{camera_init}$\to$\frm{body} and \topic{/Odometry}, starting from wherever the robot is at
launch; it drifts slowly and \pkg{dharma_localization} keeps \frm{map}$\to$\frm{odom} on top of
it. The defaults are for a Livox Mid-360 indoors.

\section{Estimators}

\begin{description}
  \item[\code{method:=fast_lio}] (default). FAST-LIO2~\cite{xu2022fastlio2}, an iterated
    error-state Kalman filter with state $(\mathbf{R}, \vect{p}, \vect{v}, \vect{b}_\omega,
    \vect{b}_a, \vect{g})$: the IMU propagates, point-to-plane residuals of each LiDAR scan
    correct. The initial gravity $\vect{g}_0$ is the mean acceleration of the first samples and
    \frm{camera_init} is the IMU frame at that time. It runs with \texttt{extrinsic\_est\_en: false}
    since the Mid-360 IMU is built in and factory calibrated (review finding \#8).
  \item[\code{method:=leg_kilo}] Leg-KILO~\cite{ou2024legkilo} through the DarmawanRobotics fork
    (headless viewer, the same \frm{odom}/\frm{body} frames as FAST-LIO, the
    \topic{/cloud_registered_body} output). It runs in LIO mode with \texttt{init\_type: 2}
    (gravity alignment: the start rotation aligns the mean acceleration with gravity). Its KILO
    mode needs per-foot contact forces that most drivers do not publish.
  \item[\code{method:=external}] another publisher of \topic{/Odometry} (the simulator); only the
    monitor runs.
\end{description}

\section{\texorpdfstring{Start gate (\code{wait_still})}{Start gate (wait_still)}}
\label{sec:wait-still}

Both estimators take gravity and the IMU biases from their first samples, so they must start
standing still. \code{wait_still} is the launch prefix of the estimator: it holds the start, and
every respawn, until a window of \param{still_sec} (2~s) of IMU samples satisfies
\begin{equation}
  \max_k \|\vect{\omega}_k\| < \omega_{\max}, \qquad
  \frac{\operatorname{std}_k \|\vect{f}_k\|}{\operatorname{mean}_k \|\vect{f}_k\|} <
  \sigma_{\mathrm{rel}}, \qquad
  \theta = \angle(\bar{\vect{f}}, \vect{z}) \le \theta_{\max},
\end{equation}
with $\omega_{\max} = 0.03$~rad/s, $\sigma_{\mathrm{rel}} = 1\,\%$ and $\theta_{\max} = 3^\circ$
(it was $8^\circ$, finding O1). The relative spread is unit-free, so an IMU reporting $g$ (the
Mid-360) or m/s$^2$ is handled alike. After \param{timeout_sec} it starts anyway with an error.

\textbf{Why the tilt matters.} If the estimator starts tilted by $\theta$, its world frame is
tilted and a horizontal distance $s$ appears as a height change and a shorter planar distance:
\begin{equation}
  \Delta z \approx s\sin\theta, \qquad \frac{\Delta s_{xy}}{s} = 1 - \cos\theta .
\end{equation}
\begin{center}
\begin{tabular}{@{}rrrr@{}}
  \toprule
  $\theta$ & $\Delta z$ per metre & distance to a 0.5~m $z$ error & planar scale error \\
  \midrule
  2° & 3.5~cm & 14.3~m & 0.06\,\% \\
  5° & 8.7~cm & 5.7~m & 0.38\,\% \\
  8° (0.1 limit) & 13.9~cm & 3.6~m & 0.97\,\% \\
  \bottomrule
\end{tabular}
\end{center}
Mapping and perception filter obstacles in a height band above the floor. With an 8° start the
floor 5~m ahead lies 0.7~m above the band of a level frame. Version 0.2.0 therefore lowers the
gate limit to 3° \emph{and} levels \frm{odom} in the localization bridge
(Section~\ref{sec:leveling}).

\section{Monitor}

\code{odometry_monitor_node} publishes \topic{odometry/status} (\code{OdometryStatus}, latched)
and \topic{odometry/restarted}.

\begin{center}
\begin{tabularx}{\textwidth}{@{}lX@{}}
  \toprule
  State & When \\
  \midrule
  \code{INITIALIZING} & no \topic{/Odometry} yet (the gate is holding the estimator) \\
  \code{LOST} & no \topic{/Odometry} for \param{stale_sec}; a jump or a restart, for
    \param{hold_sec} \\
  \code{DEGRADED} & $|z|$ above \param{max_z_drift_m} on a flat floor; \frm{odom} tilted against
    gravity by more than \param{tilt.max_deg} while still; the speed test in alarm \\
  \code{OK} & otherwise \\
  \bottomrule
\end{tabularx}
\end{center}

\textbf{Jump.} A step between consecutive poses is a jump when it is longer than
\param{jump.min_step_m} and faster than \param{jump.speed_factor} times the robot speed limit
(\param{robot.limits}).

\textbf{Restart} (review finding \#7). A fresh estimator starts at the identity pose. Being near
the origin is not enough: a robot turning in place next to its start pose gets there too. A
restart is a pose at the origin (within \param{restart.origin_radius_m} and
\param{restart.origin_yaw_deg}) after a message gap longer than \param{restart.gap_sec}
(respawn, gate) or after a translation jump back to the origin, or a stamp going back (bag loop,
simulator reset). Yaw alone never counts. On a restart the monitor publishes
\topic{odometry/restarted}: localization forgets \frm{map}$\to$\frm{odom} and levels again.

\textbf{Speed residual test.} With the robot's own odometry (\param{leg_odom_topic}), the planar
speeds are compared as a residual test with a one-sided CUSUM~\cite{page1954cusum}:
\begin{equation}
  r_k = v^{\mathrm{lio}}_k - v^{\mathrm{leg}}_k, \qquad
  \chi^2_k = \frac{r_k^2}{\sigma_{\mathrm{lio}}^2 + \sigma_{\mathrm{leg}}^2}, \qquad
  S_k = \max(0,\ S_{k-1} + \chi^2_k - \kappa),
\end{equation}
alarm while $S_k > h$. Under agreement $\chi^2_k$ follows $\chi^2(1)$ (mean 1), so
$\kappa = 3.84$ (its 95\,\% quantile) lets the sum decay, while a lasting disagreement (a
degenerate corridor, slipping feet, a stuck estimator) accumulates. The thresholds come from
the sensor standard deviations, not from per-robot speeds. $S$ is bounded at $2h$ for a fast
recovery. The test is off without leg odometry or while it is older than
\param{speed.leg_timeout_sec}.

\textbf{Tilt while still.} While the IMU is still, the gravity it measures, rotated into
\frm{odom}, gives the tilt of \frm{odom} itself; above \param{tilt.max_deg} the estimator
started tilted and $z$ will drift with distance. A robot standing on a ramp tilts its body, not
\frm{odom}, and stays \code{OK}. This assumes the IMU axes are the \topic{/Odometry} body axes
(true for FAST-LIO and Leg-KILO).

\section{Parameters}

\readonlynote

% Generated by docs/manual/tools/reference.py; do not edit.
\begin{small}
\begin{longtable}{@{}>{\raggedright}p{0.27\textwidth}>{\raggedright}p{0.14\textwidth}>{\raggedright}p{0.11\textwidth}>{\raggedright}p{0.10\textwidth}>{\raggedright\arraybackslash}p{0.26\textwidth}@{}}
\toprule
Parameter & Default & Unit & Range & Description \\
\midrule
\endhead
\bottomrule
\endfoot
\multicolumn{5}{@{}l}{\textbf{odometry\_gate (wait\_still)}} \\
\code{imu_topic}$^\dagger$ & \code{/livox/imu} &  &  & IMU topic (sensor\_msgs/Imu) \\
\code{still_sec} & 2 & s & 0.2..60 & this long without motion before the start \\
\code{max_gyro_rad_s} & 0.03 & rad/s & 0..10 & largest angular rate in the window \\
\code{max_acc_rel_std} & 0.01 &  & 0..1 & std of |acc| / mean |acc| (unit-free: m/s\textasciicircum{}2 or g) \\
\code{max_tilt_deg} & 3 & deg & 0..90 & gravity against the IMU z axis; 0: off \\
\code{min_samples} & 20 &  & 2..10000 & IMU samples needed in the window \\
\code{timeout_sec} & 60 & s & 0..3600 & start anyway (with an error); 0: wait forever \\
\multicolumn{5}{@{}l}{\textbf{odometry\_monitor\_node}} \\
\code{odom_topic}$^\dagger$ & \code{/Odometry} &  &  & estimator odometry (nav\_msgs/Odometry) \\
\code{leg_odom_topic}$^\dagger$ & \code{""} &  &  & robot leg / wheel odometry for the speed test; empty: off \\
\code{imu_topic}$^\dagger$ & \code{/livox/imu} &  &  & IMU with the axes of the odometry body; empty: no tilt \\
\code{status_period_sec}$^\dagger$ & 0.2 & s & 0.01..10 & odometry/status period \\
\code{stale_sec} & 0.5 & s & 0.05..60 & no odometry this long: LOST \\
\code{hold_sec} & 5 & s & 0..600 & a jump or restart stays LOST this long \\
\code{max_z_drift_m} & 0.5 & m & 0..100 & |z| above this is DEGRADED (flat floor); 0: off \\
\code{restart.gap_sec} & 2 & s & 0.1..60 & a message gap this long followed by the origin \\
\code{restart.origin_radius_m} & 0.05 & m & 0..1 & a fresh estimator starts this close to the origin \\
\code{restart.origin_yaw_deg} & 1 & deg & 0..180 & and with this little yaw \\
\code{jump.speed_factor} & 3 &  & 1..100 & a step faster than this times the robot speed limit \\
\code{jump.min_step_m} & 0.3 & m & 0..100 & and longer than this is a jump \\
\code{speed.window_sec} & 0.5 & s & 0.05..10 & estimator speed from the poses of this window \\
\code{speed.sigma_lio_mps} & 0.05 & m/s & 0.001..10 & std of the estimator speed \\
\code{speed.sigma_leg_mps} & 0.15 & m/s & 0.001..10 & std of the leg / wheel odometry speed \\
\code{speed.cusum_drift} & 3.84 &  & 0..1000 & kappa, subtracted from every chi\textasciicircum{}2 (chi\textasciicircum{}2(1) 95 \%: 3.84) \\
\code{speed.cusum_threshold} & 20 &  & 0.1..10000 & h: the CUSUM above this is a mismatch \\
\code{speed.leg_timeout_sec} & 0.3 & s & 0.01..10 & older leg odometry is not compared \\
\code{tilt.max_deg} & 5 & deg & 0..90 & odom z against gravity above this; 0: off \\
\code{tilt.still_sec} & 1 & s & 0.2..60 & IMU window \\
\code{tilt.max_gyro_rad_s} & 0.03 & rad/s & 0..10 & still: rate \\
\code{tilt.max_acc_rel_std} & 0.01 &  & 0..1 & still: std of |acc| / mean \\
\code{tilt.min_samples} & 20 &  & 2..10000 & IMU samples needed in the window \\
\end{longtable}
\end{small}


\section{Launch}

\begin{lstlisting}
ros2 launch dharma_odometry odometry.launch.py                  # FAST-LIO2
ros2 launch dharma_odometry odometry.launch.py method:=leg_kilo save_map:=true
ros2 launch dharma_odometry odometry.launch.py method:=external  # simulator
\end{lstlisting}

Arguments: \code{method}, \code{params_file}, \code{robot_params}, \code{site_config},
\code{use_sim_time}, \code{log_level}, \code{save_map}, \code{wait_still}, \code{respawn},
\code{leg_odom_topic}, \code{fast_lio_params_file}, \code{leg_kilo_config_file},
\code{leg_kilo_viewer}. A file argument that is set but missing fails the launch (finding
\#17). The package no longer includes the localization launch file nor records routes (moved
to \pkg{dharma_mapping}): odometry must not depend on a layer above it.
